% SPDX-License-Identifier: Apache-2.0
\section{A Reset Empties a Channel}
\label{sec:x-flush}

A channel between two units holds up to two words, and the reset
empties it. In the netlist a channel is two words and two valid bits;
the reset clears the bits, and a channel with neither bit set is
empty. The run does the same at each edge of the channel's clock while
the reset is high, whether or not anything sends or takes at that
edge, and a send or a take at that edge is dropped with the rest
(issue 729).

\code{Pair} holds a \code{Counter}, which offers its next count
whenever the channel has room, and a \code{Taker}, which takes a word
only while \code{take} is high and shows the last word it took. The
run leaves \code{take} low until the channel is full with the first
two counts, holds the reset for two cycles, and then takes a word a
cycle. What comes out starts at the counter's first count again. A
channel that kept its words through the reset would hand the taker a
one and then a zero, and the run asserts that it does not.

\begin{figure}[htbp]
\centering
\input{flush_timing.tex}
\caption{The channel between two units through a reset: \code{n}
stops at two with the channel full, \code{rst} puts it back to zero,
and \code{last} climbs from zero after it, with no word from before
the reset.}
\label{fig:x-flush}
\end{figure}

\lstinputlisting[caption={\code{ex\_flush.rs}. Run with \code{bazel run //lib/examples:ex\_flush}.},label={lst:x-flush}]{lib/examples/ex_flush.rs}

\lstinputlisting[style=ascii,caption={What \code{ex\_flush} printed when the build ran it: the channel full, the reset, then the words shown as they are taken, and the netlist, with the channel between the two units.},label={lst:x-flush-out}]{out_flush.txt}
